\SourceLeaf{140}
\Needspace{8\baselineskip}
\begin{center}
\textit{Ambige Formen.}
\end{center}
\Article{163}
Wir haben oben bereits nebenbei erwähnt, dass eine Form $F$ eine andere $F'$ sowohl eigentlich als auch uneigentlich enthalten kann. Das tritt offenbar ein, wenn sich zwischen die Formen $F,F'$ eine andere Form $G$ so einschieben lässt, dass $F$ die Form $G$ und $G$ die Form $F'$ enthält und die Form $G$ sich selbst uneigentlich äquivalent ist. Nehmen wir nämlich an, $F$ enthalte $G$ eigentlich oder uneigentlich. Da $G$ die Form $G$ uneigentlich enthält, enthält $F$ die Form $G$ auch uneigentlich beziehungsweise eigentlich, also in beiden Fällen sowohl eigentlich als auch uneigentlich (Artikel 159). Ebenso folgt daraus: Auf welche Weise auch $G$ die Form $F'$ enthalten mag, $F$ muss die Form $F'$ stets sowohl eigentlich als auch uneigentlich enthalten. --- Dass es Formen gibt, die sich selbst uneigentlich äquivalent sind, zeigt der naheliegendste Fall, in dem das mittlere Glied der Form $=0$ ist. Eine solche Form ist nämlich sich selbst entgegengesetzt (Artikel 159) und daher uneigentlich äquivalent. Allgemeiner besitzt jede Form $(a,b,c)$, in der $2b$ durch $a$ teilbar ist, diese Eigenschaft. Denn ihr ist die Form $(c,b,a)$ nach links benachbart (Artikel 160) und daher eigentlich äquivalent. Nach Artikel 159 ist aber $(c,b,a)$ der Form $(a,b,c)$ uneigentlich äquivalent, also ist $(a,b,c)$ sich selbst uneigentlich äquivalent. Solche Formen $(a,b,c)$, bei denen $2b$ durch $a$ teilbar ist, nennen wir \textit{ambig}. Es gilt also folgender Satz:

\textit{Eine Form $F$ enthält eine andere Form $F'$ sowohl eigentlich als auch uneigentlich, wenn sich eine ambige Form finden lässt, die in $F$ enthalten ist und $F'$ enthält.} Diese Aussage lässt sich aber auch umkehren, nämlich:

\Needspace{18\baselineskip}
\begin{center}
\textit{Satz über den Fall, dass eine Form zugleich eigentlich und uneigentlich in einer anderen enthalten ist.}
\end{center}
\Article{164}
\textsc{Satz.} \textit{Enthält die Form}
\[
Axx+2Bxy+Cyy\quad\mathinner{.\,.\,.\,.\,.}\quad(F)
\]
\textit{die Form}
\[
A'x'x'+2B'x'y'+C'y'y'\quad\mathinner{.\,.\,.\,.}\quad(F')
\]
\textit{sowohl eigentlich als auch uneigentlich, so lässt sich eine ambige Form finden, die in $F$ enthalten ist und die Form $F'$ enthält.}

Wir nehmen an, dass die Form $F$ in die Form $F'$ sowohl durch die Substitution
\[
x=\alpha x'+\beta y',\qquad y=\gamma x'+\delta y'
\]
\SourceLeaf{141}
als auch durch die folgende, ihr ungleichartige Substitution übergeht:
\[
x=\alpha'x'+\beta'x',\qquad y=\gamma'x'+\delta'y'
\]
Bezeichnen wir die Zahlen $\alpha\delta-\beta\gamma,\ \alpha'\delta'-\beta'\gamma'$ mit $e,e'$, so gilt $B'B'-A'C'=ee(BB-AC)=e'e'(BB-AC)$ und daher $ee=e'e'$. Da nach Voraussetzung $e,e'$ entgegengesetzte Vorzeichen haben, gilt $e=-e'$ oder $e+e'=0$. Setzt man nun in $F'$ für $x'$ den Ausdruck $\delta'x''-\beta'y''$ und entsprechend $y',-\gamma'x''+\alpha'y''$ ein, so entsteht offenbar dieselbe Form, als wenn man in $F$ schreibt:
\[
\begin{aligned}
\text{\textit{entweder} 1) für }x\quad&\alpha(\delta'x''-\beta'y'')+\beta(-\gamma'x''+\alpha'y'')\\
\text{\textit{d.\,h.}}\quad&(\alpha\delta'-\beta\gamma')x''+(\beta\alpha'-\alpha\beta')y''\\
\text{und für }y\quad&\gamma(\delta'x''-\beta'y'')+\delta(-\gamma'x''+\alpha'y'')\\
\text{\textit{d.\,h.}}\quad&(\gamma\delta'-\delta\gamma')x''+(\delta\alpha'-\gamma\beta')y''\\
\text{\textit{oder} 2) für }x\quad&\alpha'(\delta'x''-\beta'y'')+\beta'(-\gamma'x''+\alpha'y'')\quad\text{\textit{d.\,h.}}\quad e'x''\\
\text{und für }y\quad&\gamma'(\delta'x''-\beta'y'')+\delta'(-\gamma'x''+\alpha'y'')\quad\text{\textit{d.\,h.}}\quad e'y''
\end{aligned}
\]
Bezeichnen wir also die Zahlen $\alpha\delta'-\beta\gamma',\ \beta\alpha'-\alpha\beta',\ \gamma\delta'-\delta\gamma',\ \delta\alpha'-\gamma\beta'$ mit $a,b,c,d$, so wird die Form $F$ durch die beiden Substitutionen
\[
x=ax''+by'',\quad y=cx''+dy'';\qquad x=e'x'',\quad y=e'y''
\]
in dieselbe Form überführt. Daraus erhalten wir die folgenden drei Gleichungen:
\[
\begin{aligned}
Aaa+2Bac+Ccc&=Ae'e'&\quad\mathinner{.\,.\,.\,.\,.\,.\,.}\quad&[1]\\
Aab+B(ad+bc)+Ccd&=Be'e'&\quad\mathinner{.\,.\,.\,.\,.}\quad&[2]\\
Abb+2Bbd+Cdd&=Ce'e'&\quad\mathinner{.\,.\,.\,.\,.\,.\,.}\quad&[3]
\end{aligned}
\]
Aus den Werten von $a,b,c,d$ findet man aber
\[
ad-bc=ee'=-ee=-e'e'\quad\mathinner{.\,.\,.\,.\,.\,.}\quad[4]
\]
Damit ergibt sich aus $d[1]-c[2]$
\[
(Aa+Bc)(ad-bc)=(Ad-Bc)e'e'
\]
und daher
\[
A(a+d)=0
\]
Ferner ergibt sich aus $(a+d)[2]-b[1]-c[3]$

\SourceLeaf{142}

\[
(Ab+B(a+d)+Cc)(ad-bc)=(-Ab+B(a+d)-Cc)e'e'
\]

und daher

\[
B(a+d)=0
\]

Schließlich folgt aus $a[3]-b[2]$

\[
(Bb+Cd)(ad-bc)=(-Bb+Ca)e'e'
\]

und daher

\[
C(a+d)=0
\]

Da nicht alle $A,B,C$ gleich $=0$ sein können, gilt notwendig $a+d=0$ oder $a=-d$.

Aus $a[2]-b[1]$ folgt

\[
(Ba+Cc)(ad-bc)=(Ba-Ab)e'e'
\]

und somit

\[
Ab-2Ba-Cc=0\quad\mathinner{.\ .\ .\ .\ .\ .\ .\ .}\quad[5]
\]

Aus den Gleichungen $e+e'=0$, $a+d=0$ oder

\[
\alpha\delta-\beta\gamma+\alpha'\delta'-\beta'\gamma'=0,\quad\alpha\delta'-\beta\gamma'-\gamma\beta'+\delta\alpha'=0
\]

folgt $(\alpha+\alpha')(\delta+\delta')=(\beta+\beta')(\gamma+\gamma')$ oder

\[
(\alpha+\alpha'):(\gamma+\gamma')=(\beta+\beta'):(\delta+\delta')
\]

Dieses Verhältnis\HistoricalNote{*)}{Wären alle $\alpha+\alpha'$, $\gamma+\gamma'$, $\beta+\beta'$, $\delta+\delta'$ gleich $=0$, so wäre das Verhältnis unbestimmt und die Methode nicht anwendbar. Schon eine kurze Überlegung zeigt aber, dass dies mit unseren Voraussetzungen unvereinbar ist. Denn dann wäre $\alpha\delta-\beta\gamma=\alpha'\delta'-\beta'\gamma'$, das heißt $e=e'$, und daher wegen $e=-e'$ auch $e=e'=0$. Damit würde auch $B'B'-A'C'$, also die Determinante der Form $F'$, gleich $=0$. Solche Formen haben wir vollständig ausgeschlossen.} werde in kleinsten Zahlen als $m:n$ geschrieben, sodass $m,n$ teilerfremd sind. Man wähle $\mu,\nu$ so, dass $\mu m+\nu n=1$ gilt. Ferner sei $r$ der größte gemeinsame Teiler der Zahlen $a,b,c$. Sein Quadrat teilt daher $aa+bc$ oder $bc-ad$ oder $ee$, also teilt $r$ auch $e$. Nehmen wir bei diesen Festlegungen an, dass die Form $F$ durch die Substitution

\[
x=mt+\frac{\nu e}{r}u,\qquad y=nt-\frac{\mu e}{r}u
\]

in die Form $Mtt+2Ntu+Puu$, $(G)$ übergeht, so ist diese ambig und enthält die Form $F'$.

\SourceLeaf{143}

\emph{Beweis.} I. Um zu zeigen, dass die Form $G$ ambig ist, beweisen wir

\[
M(b\mu\mu-2a\mu\nu-c\nu\nu)=2Nr
\]

Da $r$ die Zahlen $a,b,c$ teilt, ist $\frac{1}{r}(b\mu\mu-2a\mu\nu-c\nu\nu)$ eine ganze Zahl und daher $2N$ ein Vielfaches von $M$. Es gilt aber

\[
\begin{gathered}M=Amm+2Bmn+Cnn,\\Nr=(Am\nu-B(m\mu-n\nu)-Cn\mu)e\quad\mathinner{.\ .}\quad[6]\end{gathered}
\]

Ferner bestätigt man durch Ausmultiplizieren leicht, dass

\[
\begin{aligned}2e+2a&=e-e'+a-d\\&=(\alpha-\alpha')(\delta+\delta')-(\beta-\beta')(\gamma+\gamma')\end{aligned}
\]

\[
2b=(\alpha+\alpha')(\beta-\beta')-(\alpha-\alpha')(\beta+\beta')
\]

gilt. Wegen $m(\gamma+\gamma')=n(\alpha+\alpha')$, $m(\delta+\delta')=n(\beta+\beta')$ ergibt sich daraus

\[
m(2e+2a)=-2nb\quad\text{oder}
\]

\[
me+ma+nb=0\quad\mathinner{.\ .\ .\ .\ .\ .\ .\ .\ .}\quad[7]
\]

Ebenso gilt

\[
\begin{aligned}2e-2a&=e-e'-a+d\\&=(\alpha+\alpha')(\delta-\delta')-(\beta+\beta')(\gamma-\gamma')\end{aligned}
\]

\[
2c=(\gamma-\gamma')(\delta+\delta')-(\gamma+\gamma')(\delta-\delta')
\]

und daher $n(2e-2a)=-2mc$ oder

\[
ne-na+mc=0\quad\mathinner{.\ .\ .\ .\ .\ .\ .\ .\ .}\quad[8]
\]

Addiert man nun zu $mm(b\mu\mu-2a\mu\nu-c\nu\nu)$ den Ausdruck

\[
\begin{gathered}(1-m\mu-n\nu)(m\nu(e-a)+(m\mu+1)b)\\+(me+ma+nb)(m\mu\nu+\nu)+(ne-na+mc)m\nu\nu\end{gathered}
\]

der offenbar $=0$ ist, weil

\[
1-\mu m-\nu u=0,\quad me+ma+nb=0,\quad ne-na+mc=0
\]

gilt, so ergibt sich nach Ausmultiplizieren der Produkte und Weglassen der sich aufhebenden Glieder $2m\nu e+b$. Also gilt

\[
mm(b\mu\mu-2a\mu\nu-c\nu\nu)=2m\nu e+b\quad\mathinner{.\ .\ .\ .\ .}\quad[9]
\]

Addiert man ebenso zu $mn(b\mu\mu-2a\mu\nu-c\nu\nu)$ die Ausdrücke

\SourceLeaf{144}

\[
\begin{gathered}(1-m\mu-n\nu)((n\nu-m\mu)e-(1+m\mu+n\nu)a)\\-(me+ma+nb)m\mu\mu+(ne-na+mc)n\nu\nu\end{gathered}
\]

so erhält man

\[
mn(b\mu\mu-2a\mu\nu-c\nu\nu)=(n\nu-m\mu)e-a\quad .\ .\ .\ .\ [10]
\]

Addiert man schließlich zu $nn(b\mu\mu-2a\mu\nu-c\nu\nu)$ die Ausdrücke

\[
\begin{gathered}(m\mu+n\nu-1)(n\mu(e+a)+(n\nu+1)c)\\-(me+ma+nb)n\mu\mu-(ne-na+mc)(n\mu\nu+\mu)\end{gathered}
\]

so ergibt sich

\[
nn(b\mu\mu-2a\mu\nu-c\nu\nu)=-2n\mu e-c\quad .\ .\ .\ .\ [11]
\]

Aus $9$, $10$, $11$ leitet man nun ab:

\[
\begin{gathered}(Amm+2Bmn+Cnn)(b\mu\mu-2a\mu\nu-c\nu\nu)\\=2e(Am\nu+B(n\nu-m\mu)-Cn\mu)+Ab-2Ba-Cc\end{gathered}
\]

oder wegen $[6],$

\[
M(b\mu\mu-2a\mu\nu-c\nu\nu)=2Nr.\quad \text{W.\,z.\,b.\,w.}
\]

II. Um zu beweisen, dass die Form $G$ die Form $F'$ enthält, zeigen wir \emph{erstens}, dass $G$ in $F'$ übergeht, wenn man

\[
\begin{gathered}t=(\mu\alpha+\nu\gamma)x'+(\mu\beta+\nu\delta)y',\\u=\frac re(n\alpha-m\gamma)x'+\frac re(n\beta-m\delta)y'\ .\ .\ .\ (S)\end{gathered}
\]

setzt, und \emph{zweitens}, dass $\frac{r}{e}(n\alpha-m\gamma)$, $\frac{r}{e}(n\beta-m\delta)$ ganze Zahlen sind.

1. Da $F$ in $G$ übergeht, wenn man

\[
x=mt+\frac{\nu e}{r}u,\quad y=nt-\frac{\mu e}{r}u
\]

setzt, wird die Form $G$ durch die Substitution $(S)$ in dieselbe Form überführt, in die $F$ übergeht, wenn man

\[
\begin{aligned}x={}&m((\mu\alpha+\nu\gamma)x'+(\mu\beta+\nu\delta)y')\\&+\nu((n\alpha-m\gamma)x'+(n\beta-m\delta)y')\end{aligned}
\]

\emph{das heißt} $=\alpha(m\mu+n\nu)x'+\beta(m\mu+n\nu)y'$ oder $=\alpha x'+\beta y'$,

und

\[
\begin{aligned}y={}&n((\mu\alpha+\nu\gamma)x'+(\mu\beta+\nu\delta)y')\\&-\mu((n\alpha-m\gamma)x'+(n\beta-m\delta)y')\end{aligned}
\]

\emph{das heißt} $=\gamma(n\nu+m\mu)x'+\delta(n\nu+m\mu)y'$ oder $=\gamma x'+\delta y'$ setzt.

\SourceLeaf{145}

Durch diese Substitution geht aber $F$ in $F'$ über. Also geht durch die Substitution $(S)$ auch $G$ in $F'$ über.

2. Aus den Werten von $e,b,d$ findet man $\alpha'e+\gamma b-\alpha d=0$ oder wegen $d=-a$ die Gleichung $n\alpha'e+n\alpha a+n\gamma b=0$. Daraus folgt nach $[7]$ die Gleichung $n\alpha'e+n\alpha a=m\gamma e+m\gamma a$ oder

\[
(n\alpha-m\gamma)a=(m\gamma-n\alpha')e\quad .\ .\ .\ .\ .\ .\ .\ [12]
\]

Ferner gilt $\alpha nb=-\alpha m(e+a)$, $\gamma mb=-m(\alpha'e+\alpha a)$, und daher

\[
(n\alpha-m\gamma)b=(\alpha'-\alpha)me\quad .\ .\ .\ .\ .\ .\ .\ [13]
\]

Schließlich gilt $\gamma'e-\gamma a+\alpha c=0$. Multipliziert man mit $n$ und setzt für $na$ seinen Wert aus $[8]$ ein, so folgt

\[
(n\alpha-m\gamma)c=(\gamma-\gamma')ne\quad .\ .\ .\ .\ .\ .\ .\ [14]
\]

Ebenso ergibt sich $\beta'e+\delta b-\beta d=0$ oder $n\beta'e+n\delta b+n\beta a=0$ und damit nach $[7],$ die Gleichung $n\beta'e+n\beta a=m\delta e+m\delta a$ oder

\[
(n\beta-m\delta)a=(m\delta-n\beta')e\quad .\ .\ .\ .\ .\ .\ .\ [15]
\]

Ferner gilt $\beta nb=-\beta m(e+a)$, $\delta mb=-m(\beta'e+\beta a)$, und daher

\[
(n\beta-m\delta)b=(\beta'-\beta)me\quad .\ .\ .\ .\ .\ .\ .\ [16]
\]

Schließlich gilt $\delta e-\delta a+\beta c=0$. Multipliziert man mit $n$ und setzt für $na$ seinen Wert aus $[8]$ ein, so erhält man

\[
(n\beta-m\delta)c=(\delta-\delta')ne\quad .\ .\ .\ .\ .\ .\ .\ [17]
\]

Da nun der größte gemeinsame Teiler der Zahlen $a,b,c$ gleich $r$ ist, lassen sich ganze Zahlen $\mathfrak{A},\mathfrak{B},\mathfrak{C}$ so wählen, dass

\[
\mathfrak{A}a+\mathfrak{B}b+\mathfrak{C}c=r
\]

gilt. Damit folgt aus $12$, $13$, $14$ und $15$, $16$, $17$

\[
\begin{aligned}&\mathfrak A(m\gamma-n\alpha')+\mathfrak B(\alpha'-\alpha)m+\mathfrak C(\gamma-\gamma')n\\&\qquad=\frac re(n\alpha-m\gamma)\\&\mathfrak A(m\delta-n\beta')+\mathfrak B(\beta'-\beta)m+\mathfrak C(\delta-\delta')n\\&\qquad=\frac re(n\beta-m\delta)\end{aligned}
\]

also sind $\frac{r}{e}(n\alpha-n\gamma),\quad\frac{r}{e}(n\beta-m\delta)$ ganze Zahlen. $\text{W.\,z.\,b.\,w.}$

\SourceLeaf{146}

\Article{165}

\emph{Beispiel.} Die Form $3xx+14xy-4yy$ wird in die Form $-12x'x'-18x'y'+39y'y'$ überführt, sowohl eigentlich durch die Substitution

\[
x=4x'+11y',\quad y=-x'-2y'
\]

als auch uneigentlich durch die Substitution

\[
x=-74x'+89y',\quad y=15x'-18y'
\]

Hier sind also $\alpha+\alpha'$, $\beta+\beta'$, $\gamma+\gamma'$, $\delta+\delta'$ der Reihe nach $-70$, $100$, $14$, $-20$, und es gilt $-70:14=100:-20=5:-1$. Wir setzen daher $m=5$, $n=-1$, $\mu=0$, $\nu=-1$. Die Zahlen $a,b,c$ ergeben sich als $-237$, $-1170$, $48$, mit dem größten gemeinsamen Teiler $=3=r$. Schließlich gilt $e=3$. Damit lautet die Transformation $(S)$: $x=5t-u$, $y=-t$. Durch sie geht die Form $(3,7,-4)$ in die ambige Form $tt-16tu+3uu$ über.

\Needspace{14\baselineskip}
Sind die Formen $F,F'$ äquivalent, so ist die Form $G$, die in $F$ enthalten ist, auch in $F'$ enthalten. Da sie aber dieselbe Form auch enthält, ist sie ihr und damit auch der Form $F$ äquivalent. In diesem Fall lässt sich der Satz also so aussprechen:

\emph{Sind $F,F'$ sowohl eigentlich als auch uneigentlich äquivalent, so lässt sich eine beiden äquivalente ambige Form finden.} --- In diesem Fall ist $e=\pm1$, und daher ist auch $r.$, das $e$ teilt, gleich $=1$.

Dies möge über die Transformation der Formen im Allgemeinen genügen. Wir gehen nun zur Betrachtung der \emph{Darstellungen} über.
